\chapter{Identificación de Requisitos}
\label{cap:requisitos}

El trabajo tiene dos núcleos y por tanto dos conjuntos de requisitos de naturaleza
distinta: los que debe cumplir la aplicación para ser útil en un aula, y los que debe
cumplir el procedimiento experimental para que sus conclusiones sean válidas. Los segundos
suelen quedar implícitos; aquí se enuncian porque varias decisiones del trabajo se derivan
directamente de ellos.

\section{Requisitos del sistema de subtitulado}

\subsection{Funcionales}

Los requisitos funcionales, recogidos en la tabla~\ref{tab:rf}, describen lo que el sistema
debe hacer desde el punto de vista de quien lo usa: el docente que emite y el alumnado que
lee.

\begin{table}[h]
\centering
\small
\begin{tabular}{|c|p{11cm}|}
\hline
\textbf{Id} & \textbf{Requisito} \\
\hline
RF1 & El sistema transcribe en directo la voz del docente y muestra el resultado como
      subtítulos. \\
\hline
RF2 & Varios alumnos pueden visualizar simultáneamente los subtítulos desde sus propios
      dispositivos. \\
\hline
RF3 & El alumnado no necesita instalar nada ni conceder permisos: basta con abrir una
      dirección en el navegador. \\
\hline
RF4 & El docente puede definir, antes de comenzar, la terminología de la sesión, que se
      resalta visualmente en los subtítulos. \\
\hline
RF5 & Un alumno que se incorpora una vez iniciada la clase recupera lo transcrito hasta
      ese momento. \\
\hline
RF6 & Cada alumno ajusta el tamaño del texto en su propio dispositivo. \\
\hline
RF7 & El sistema informa de la latencia con que está operando. \\
\hline
\end{tabular}
\caption{Requisitos funcionales del sistema de subtitulado.}
\label{tab:rf}
\end{table}

El requisito RF6 no es un detalle de comodidad. El sistema se dirige a personas con una
discapacidad sensorial, y no cabe suponer que quien lo desarrolla pueda fijar un tamaño de
texto adecuado para todas ellas ni para todas las distancias de lectura de un aula.

RF3 tiene una justificación menos evidente y más importante de lo que aparenta. La
población que se beneficia de un subtítulo de aula es mayor que el colectivo que lo motiva:
incluye al alumnado que aún no domina la lengua de la clase, con evidencia empírica detrás
(sección~\ref{sec:para-quien}). Un recurso que hubiera que solicitar, instalar o activar
excluiría por diseño a casi toda esa población, y obligaría a identificarse a quien
prefiriera no hacerlo. De ahí que el requisito no sea «que sea fácil de usar» sino que no
haya nada que hacer: abrir una dirección.

\subsection{No funcionales}

Los de la tabla~\ref{tab:rnf} no describen funciones sino condiciones bajo las cuales las
anteriores son útiles: de poco sirve transcribir bien si el subtítulo llega tarde, si el
sistema necesita una unidad de cálculo por alumno o si el audio del aula sale del centro.

\begin{table}[h]
\centering
\small
\begin{tabular}{|c|p{11cm}|}
\hline
\textbf{Id} & \textbf{Requisito} \\
\hline
RNF1 & La latencia entre lo dicho y lo mostrado debe permitir seguir la clase. El umbral
       concreto se determina experimentalmente y se reporta con su valor típico y su peor
       caso. \\
\hline
RNF2 & El coste de cálculo no debe crecer con el número de alumnos conectados. \\
\hline
RNF3 & El sistema opera exclusivamente en la red local del centro. \\
\hline
RNF4 & El puesto de emisión está protegido frente a uso no autorizado. \\
\hline
RNF5 & El sistema funciona con equipamiento asequible para un centro educativo. \\
\hline
RNF6 & El requisito de captación de audio debe expresarse en términos verificables, no
       como recomendación genérica. \\
\hline
\end{tabular}
\caption{Requisitos no funcionales del sistema de subtitulado.}
\label{tab:rnf}
\end{table}

RNF2 es el que determina la arquitectura. Si cada dispositivo del aula ejecutara el modelo,
el sistema exigiría una unidad de cálculo por alumno y dejaría de ser desplegable. De ahí
que todo el reconocimiento resida en el servidor y los clientes se limiten a visualizar.

RNF3 se deriva de la naturaleza del material: el sistema difunde audio de aula transcrito,
con voces identificables y frecuentemente de menores de edad. La restricción no puede
delegarse en la configuración de red del centro, porque un reenvío de puertos mal
configurado la anularía sin que nadie lo advirtiera; debe imponerla la propia aplicación.

RNF6 responde a un problema recurrente en los requisitos de este tipo de sistemas. Indicar
que se necesita ``un micrófono adecuado'' no es verificable. La caracterización
experimental frente al ruido permite sustituirlo por una condición comprobable sobre la
relación señal-ruido, y con ella justificar la elección del equipo de captación.

\subsection{Legales y de protección de datos}
\label{sec:requisitos-legales}

Un sistema que capta la voz de un aula y la difunde por una red trata datos personales, y
frecuentemente de menores. Los requisitos que se derivan de ello no son una formalidad
añadida al final: dos de los recogidos en la tabla~\ref{tab:rl} determinaron decisiones
de arquitectura.

\begin{table}[h]
\centering
\small
\begin{tabular}{|c|p{11cm}|}
\hline
\textbf{Id} & \textbf{Requisito} \\
\hline
RL1 & La voz captada no abandona la red del centro ni se envía a ningún servicio de
      terceros. El reconocimiento se ejecuta en un equipo del propio centro. \\
\hline
RL2 & El sistema no almacena el audio. La transcripción vive en memoria mientras dura la
      sesión y se descarta al terminarla. \\
\hline
RL3 & La restricción de alcance la impone la aplicación, no la configuración de red del
      centro, de modo que sea verificable y no dependa de una suposición. \\
\hline
RL4 & El material empleado en la comparativa procede de corpus con licencia de uso para
      investigación; no se grabaron clases reales, y por tanto no se recabaron
      consentimientos ni se trataron datos de alumnado identificable. \\
\hline
\end{tabular}
\caption{Requisitos legales y de protección de datos.}
\label{tab:rl}
\end{table}

RL1 es la razón de que el reconocimiento se ejecute en local en lugar de delegarse en un
servicio comercial en la nube, que habría sido más sencillo y probablemente más preciso.
Enviar audio de aula a un tercero exigiría una base jurídica y un contrato de encargo de
tratamiento \cite{ue2016rgpd} que quedan fuera de lo que un centro puede resolver por su
cuenta para una asignatura. Ejecutarlo en local convierte un problema de cumplimiento en un
problema de ingeniería, que es la razón de fondo de la restricción de recursos que atraviesa
todo el trabajo.

RL4 delimita el alcance ético de la parte experimental: al no grabarse material propio, no
hubo tratamiento de datos personales que requiriera consentimiento informado ni evaluación
ética. La continuación prevista en el capítulo~\ref{cap:conclusiones} es repetir la
comparativa sobre clases reales; si llega a ejecutarse, ese punto deja de ser trivial y debe
resolverse antes de grabar, no después.

\section{Requisitos del procedimiento experimental}

Estos requisitos, listados en la tabla~\ref{tab:re}, no proceden del usuario sino de la
naturaleza de la pregunta que el trabajo pretende responder. Se enuncian porque su
incumplimiento no produce un fallo visible: produce cifras verosímiles y falsas.

\begin{table}[h]
\centering
\small
\begin{tabular}{|c|p{11cm}|}
\hline
\textbf{Id} & \textbf{Requisito} \\
\hline
RE1 & Toda cifra publicada debe poder regenerarse a partir del código y los datos
       registrados junto a ella. \\
\hline
RE2 & La identidad de un resultado incluye su configuración completa, no solo el modelo y
       el corpus. \\
\hline
RE3 & El tamaño de muestra debe justificarse antes de interpretar diferencias, y las
       mediciones que no lo alcancen deben marcarse como no concluyentes. \\
\hline
RE4 & Las comparaciones entre condiciones deben ser pareadas y compartir configuración. \\
\hline
RE5 & Los datos empleados para adaptar el sistema deben ser disjuntos de los empleados
       para evaluarlo. \\
\hline
RE6 & La calidad de las referencias debe auditarse antes de adoptarlas. \\
\hline
RE7 & Ninguna cifra de la memoria se transcribe manualmente: las tablas y figuras se
       generan desde los datos de ejecución. \\
\hline
\end{tabular}
\caption{Requisitos del procedimiento experimental.}
\label{tab:re}
\end{table}

RE2 merece justificación porque resulta menos evidente que los demás. Durante el trabajo se
comprobó que dos ejecuciones del mismo modelo sobre el mismo corpus, diferenciadas
únicamente por la estrategia de decodificación, arrojaban tasas de error que diferían en
más de doce puntos. Un resultado identificado solo por modelo y corpus habría permitido
comparar cifras no comparables sin que nada lo delatara.

RE6 se incorporó tras comprobar que una parte sustancial del error atribuido al modelo en
determinados corpus correspondía en realidad a defectos de la transcripción de referencia.

RE7 es una consecuencia práctica de RE1: una cifra copiada a mano deja de estar vinculada a
la ejecución que la produjo en el momento en que se copia.

\section{Trazabilidad de los requisitos}
\label{sec:trazabilidad-requisitos}

Enunciar requisitos sin cerrarlos deja al lector sin forma de comprobar si se cumplieron.
La tabla~\ref{tab:trazabilidad} los vincula con el mecanismo concreto que los satisface y
con la evidencia que lo demuestra; el capítulo~\ref{cap:desarrollo} desarrolla cada uno y el
capítulo~\ref{cap:conclusiones} verifica los objetivos que de ellos se derivan.

\begin{table}[h]
\centering
\footnotesize
\setlength{\tabcolsep}{4pt}
\begin{tabular}{|c|p{5.3cm}|p{5.3cm}|}
\hline
\textbf{Id} & \textbf{Dónde se resuelve} & \textbf{Cómo se comprueba} \\
\hline
RF1 & Segmentador en vivo y servicio de reconocimiento & Demostración funcional; sesión con
      micrófono real \\
\hline
RF2 & Sesión de subtitulado compartida, difusión por websocket & Pruebas automáticas de la
      sesión; varios clientes simultáneos \\
\hline
RF3 & Vista de alumnado sin captura ni permisos & Abrir la dirección en un navegador sin
      instalar nada \\
\hline
RF4 & Decorador de resaltado sobre el motor de reconocimiento & Pruebas del glosario,
      incluida la protección de la eñe \\
\hline
RF5 & Histórico retenido en la sesión y enviado al conectar & Prueba automática de
      incorporación tardía \\
\hline
RF6 & Control de tamaño de texto en la vista del alumnado & Verificación manual en la
      interfaz \\
\hline
RF7 & Instrumentación de latencia de extremo a extremo en el navegador & Panel de la vista
      de emisión (sección~\ref{sec:latencia}) \\
\hline
RNF1 & Segmentación por pausas con tope de duración & exp-100 y exp-102; medición del
       circuito completo \\
\hline
RNF2 & Reconocimiento único en el servidor para toda el aula & Arquitectura; coste
       independiente del número de clientes \\
\hline
RNF3 & Filtro de red privada en la propia aplicación & Pruebas automáticas del filtro \\
\hline
RNF4 & Clave compartida en el puesto de emisión & Pruebas automáticas de acceso \\
\hline
RNF5 & Modelo de tamaño medio sobre una GPU de consumo & exp-000 y exp-101 \\
\hline
RNF6 & Curva de degradación frente al ruido & exp-104: requisito expresado en dB \\
\hline
RL1--RL3 & Ejecución local, sin persistencia de audio, filtro de red privada & Pruebas
       automáticas y revisión de la configuración de despliegue \\
\hline
RE1--RE7 & Registro de procedencia y generación automática de tablas y figuras &
       Apéndice~\ref{ap:reproducibilidad} \\
\hline
\end{tabular}
\caption{Trazabilidad de requisitos: dónde se resuelve cada uno y con qué evidencia.}
\label{tab:trazabilidad}
\end{table}
